\chapter{The catenoid as an eternal stationary flow}\label{sec:catalogue-stationary-catenoid}
\begin{example}[Stationary catenoid flow]\label{ex:catalogue-stationary-catenoid}
For $a>0$ the explicit immersion
\[
 X(\theta,s,t)=a(\cosh s\cos\theta,\cosh s\sin\theta,s),
 \quad \theta\in\mathbb R/2\pi\mathbb Z,\quad s\in\mathbb R,\quad t\in\mathbb R,
\]
is independent of time and has $\vec H=0$. It is therefore an eternal stationary mean curvature flow in $\mathbb R^3$.

Every time slice is a complete properly embedded noncompact rotational annulus, with two ends and a waist circle of radius $a$. Its metric and curvature are
\[
 g=a^2\cosh^2s(d\theta^2+ds^2),\qquad
 K=-\frac1{a^2\cosh^4s},\qquad
 \int K\,dA=-4\pi.
\]
The metric follows by differentiation, and the conformal curvature formula gives $K$ and its integral. The complete catenoid is unstable for compactly supported area variations: for a test function $f(s)$ independent of $\theta$, its second variation is
\[
 Q(f,f)=2\pi\int_{\mathbb R}\bigl((f')^2-2\operatorname{sech}^2s\,f^2\bigr)\,ds.
\]
Taking $f$ equal to one on $[-R,R]$, zero outside $[-2R,2R]$, and $|f'|\leq C/R$ makes this negative for large $R$. Each open half is a minimal graph over $r>a$, with $z=\pm a\operatorname{arcosh}(r/a)$, so it is stable for compactly supported variations in that half. These are minimal-area stability assertions, not a classification of dynamical stability under MCF. The supplied source also records the catenoid in \cite{CatMan11}, Remark 4.4.7, printed p.~105, supplied PDF p.~116. The parametrization is exact and supports interactive rotation, scale selection, and a truncated height range with the omitted infinite ends stated.
Chopp uses the bounded piece of this classical surface as a numerical
benchmark in \cite{CatChoppMinimal1993}, Section~4.1, pp.~86--87.
In axial coordinates its exact profile is
$r(z)=a\cosh(z/a)$, $-b\leq z\leq b$, with the two fixed boundary
circles of radius $R=a\cosh(b/a)$. The test uses $a=0.4$ and
$b\approx0.277259$, giving $R\approx0.5$. This is a fixed-boundary
minimal-surface computation, not a new complete catenoid or a closed
formula for the transient numerical evolution. The solver's artificial
time and curvature conventions should not be identified with the
normalized MCF time used elsewhere in this collection.
\end{example}

\chapter{Mean curvature evolution of the Whitney sphere}\label{sec:catalogue-whitney-sphere}
\begin{example}[A collapsing Lagrangian immersed sphere]\label{ex:catalogue-whitney-sphere}
For $n\geq2$, the initial Whitney immersion in $\mathbb R^{2n}=\mathbb C^n$ is explicitly
\[
 F_0(x_1,\ldots,x_{n+1})
 =\frac{(x_1,x_1x_{n+1},\ldots,x_n,x_nx_{n+1})}{1+x_{n+1}^2},
 \qquad (x_1,\ldots,x_{n+1})\in S^n.
\]
It is a compact Lagrangian immersion, of dimension and codimension $n$, with one transverse double image point at the origin coming from the two poles. It is complete and proper as an immersion of a compact domain, and it is not embedded. Its equivariance reduces the geometry to the planar initial profile
\[
 \gamma_0(u)=\left(\frac{\sin u}{1+\cos^2u},
 \frac{\sin u\cos u}{1+\cos^2u}\right),
 \qquad u\in\mathbb R/2\pi\mathbb Z.
\]
For $n\geq2$ the profile evolution is the equivariant equation
\[
 (\partial_t\gamma)^\perp
 =\vec k-(n-1)\frac{\gamma^\perp}{|\gamma|^2},
\]
interpreted by the smooth equivariant immersion at the double point. It is not ordinary planar curve shortening.

The maximal smooth mean curvature flow $F_t$, $t\in[0,T)$, remains Lagrangian and equivariant. It collapses to the origin in finite time, and its tangent flow there is a Lagrangian plane with multiplicity two. The Whitney sphere is not a self-similar solution. These statements are the special case stated in \cite{CatViana18}, pp.~1--3, especially the result on p.~2 and Theorem 1.3 on p.~3. The formulas above specify the initial geometry and its reduced equation, not a closed expression for $F_t$ at positive time. The illustration specializes to $n=2$. Write $\gamma_t=a+ib$. A single right-hand lobe, from one preimage of the origin to the other, generates the immersed two-sphere by
\[
 F_t(u,v)=(\gamma_t(u)\cos v,\gamma_t(u)\sin v)\in\mathbb C^2,
 \qquad v\in\mathbb R/2\pi\mathbb Z.
\]
The circles at the two endpoints collapse to the two distinct poles of the domain, which have the same image at the origin. Using both antipodal lobes would cover the same sphere twice.

The time sequence pairs this surface, linearly projected from $\mathbb R^4$ to $\mathbb R^3$, with its planar profile. At $t=0$ the geometry is exact; the positive-time frames are numerical approximations of the equivariant PDE for $n=2$, with $\partial_tF=\vec H$ and the initial scale fixed above. They show $0\leq t\leq0.09$ on a common spatial scale, with no normalization by an estimated extinction time. The computation stops before the singularity: $0.09$ is not asserted to equal $T$. Spatial mesh refinement and tighter time-integration tolerances are compared, but these checks are not certified error bounds. The profile narrows and contracts rather than shrinking by a single scale factor. Its double point remains at the origin; further overlaps in the surface projection need not be intersections in $\mathbb R^4$. The multiplicity-two tangent plane is a parabolic blow-up limit, not an unscaled final surface displayed by the animation.
\end{example}

\chapter{Rotational degenerate neckpinch solutions}\label{sec:catalogue-degenerate-neckpinches}
\begin{example}[Angenent--Vel\'azquez type-II singularities]\label{ex:catalogue-degenerate-neckpinches}
Let $n\geq2$, let $m\geq3$ be odd, and let $K>0$. Choose a compact rotational target hypersurface in $\mathbb R^{n+1}$, smooth away from a cusp at the origin, whose radius satisfies $u_*(x)\sim\sqrt K\,x^{m/2}$ as $x\downarrow0$ and whose other tip closes smoothly. Angenent and Vel\'azquez construct smooth compact rotational mean curvature flows
\[
 M_t=\{(x,u(x,t)\omega):a(t)\leq x\leq b(t),\ \omega\in S^{n-1}\},
 \qquad t_0<t<T,
\]
with smooth caps and positive radius between them. Each regular slice is an embedded sphere, hence complete and proper. The constructed flows remain close to the chosen target outside a small neighborhood of the origin as the construction interval is shortened. They are not supplied as ancient solutions or as solitons, and no closed formula for $u(x,t)$ is given.

Set $\xi_*=(2(n-1)/K)^{1/m}$. At the left tip,
\[
 a(t)=-(\xi_*+o(1))(T-t)^{1/m},\qquad
 |\vec H|_{\rm tip}(t)=\left(\frac{\xi_*}{m}+o(1)\right)(T-t)^{1/m-1}.
\]
Here the curvature coefficient is for the trace mean curvature, equivalently the magnitude $|\vec H|$ at this convex tip. Rotational smoothness makes the $n$ principal curvatures equal at the tip, so $|A|_{\rm tip}=|\vec H|_{\rm tip}/\sqrt n$. The maximum of $|A|$ is attained there; its blow-up rate is type II. The parabolic rescalings at the origin converge smoothly on compact sets to $S^{n-1}_{\sqrt{2(n-1)}}\times\mathbb R$, whereas rescaling about the tip by $A_{\max}(t)=\max_{M_t}|A|$ yields the rotational translating bowl with tip $|A|=1$. This last normalization has translation speed $\sqrt n$; rescaling instead by $|\vec H|_{\rm tip}$ gives the speed-one bowl.

Under the intermediate scaling $x\mapsto x/(T-t)^{1/m}$ and $r\mapsto r/(T-t)^{1/2}$, the limiting profile is
\[
 r=\sqrt{2(n-1)+Kx^m},\qquad x\geq-\xi_*.
\]
This last expression is an asymptotic limiting profile, not an exact formula for the evolving compact hypersurface. Source: \cite{CatAV97}, Introduction, items 1--5, printed pp.~16--17, supplied PDF pp.~2--3; the construction is justified in Sections 3--7. The even-$m$ modifications and the formal $m=2$ discussion are not included in this entry. The illustration selects $n=2$, $m=3$, $K=1$, and writes $s=T-t$. Its time sequence uses one spatial scale to show the left tip approaching the origin while the outer cap persists; the final cusp target is not a smooth time slice. For the profile near the origin, item 5 of the Introduction gives
\[
 \frac{u(y\sqrt{s},T-s)}{\sqrt{s}}
 =\sqrt{2}+\frac{\sqrt{s}}{2\sqrt{2}}(y^3-6y)+o(\sqrt{s})
\]
on bounded $y$ intervals. This leading Hermite correction explains the small bulge and neck in the cylindrical region. The rotatable compact slices blend the asymptotic profiles with a numerically integrated bowl cap and a chosen outer cap: they are illustrative schematics, not a computed solution of the compact MCF initial-value problem, and their displayed times carry no certified finite-time accuracy. Separate views compare the meridians and the three blow-up limits; the intermediate view uses unequal axial and radial scaling. The construction does not classify all type-II singularities or impose global convexity on every allowed target.
\end{example}

\chapter{The Angenent oval as an explicit ancient flow}\label{sec:catalogue-angenent-oval}
\begin{example}[The ancient paperclip curve]\label{ex:catalogue-angenent-oval}
With the speed-one curve-shortening convention, the exact implicit equation
\[
 \Gamma_t=\{(x,y)\in\mathbb R^2:|x|<\pi/2,\ e^t\cosh y=\cos x\},
 \qquad -\infty<t<0,
\]
defines the Angenent oval. Equivalently, its two graph branches are
\[
 y=\pm\operatorname{arcosh}(e^{-t}\cos x),
 \qquad |x|\leq\arccos(e^t).
\]
Their endpoints join smoothly as a curve, even though those graph charts cease to be regular there. This is an exact implicit formula, rather than a single global graph parametrization. It is recorded in \cite{BV11}, Section 3.1, equation (14), p.~7.

Every time slice is a smooth compact embedded strictly convex curve, hence complete and proper, with topology $S^1$ and reflection symmetries across both coordinate axes. It shrinks to the origin at time zero but changes shape and is not a soliton. Its enclosed area is $-2\pi t$, and
\[
 \kappa_{\max}=(1-e^{2t})^{-1/2},\qquad
 \kappa_{\min}=e^t(1-e^{2t})^{-1/2}.
\]
The rescaling $(-t)^{-1/2}\Gamma_t$ converges smoothly to the radius-$\sqrt2$ circle at extinction. As $t\to-\infty$, translating its upper and lower branches to keep their tips fixed gives the two oppositely directed grim-reaper profiles, locally in $|x|<\pi/2$. These limits can be read from the implicit equation: near extinction its quadratic expansion gives $x^2+y^2\sim-2t$, and subtracting the upper tip height gives the limiting graph $\log\cos x$ as $t\to-\infty$. The time slider plots the exact implicit curve for $-5\leq t\leq-0.1$ on a fixed spatial scale.
\end{example}

\chapter{Ancient torus curves in higher codimension}\label{sec:catalogue-ancient-torus-curves}
\begin{example}[Compact ancient torus curves]\label{ex:catalogue-ancient-torus-curves}
Fix $m\geq2$ and integers $1=k_1<k_2<\cdots<k_m$. For $t<0$, let $r(t)>0$ be the unique solution of
\[
 \sum_{j=1}^m r(t)^{2k_j^2}=-2t.
\]
Then the exact parametrization
\[
 \gamma_t(\theta)=
 \bigl(r(t)^{k_1^2}\cos(k_1\theta),r(t)^{k_1^2}\sin(k_1\theta),
 \ldots,r(t)^{k_m^2}\cos(k_m\theta),r(t)^{k_m^2}\sin(k_m\theta)\bigr),
 \qquad \theta\in\mathbb R/2\pi\mathbb Z,
\]
is an ancient curve-shortening flow in $\mathbb R^{2m}$. Its scale is specified by an exact implicit equation; no elementary inverse in time is required. Equivalently,
\[
 r'=-\frac{r}{\sum_j k_j^2r^{2k_j^2}}.
\]
Each time slice is a compact smooth embedded circle, hence complete and proper, with no ends and codimension $2m-1$. Embeddedness follows directly from $k_1=1$, because the first two coordinates already determine $\theta$. The name refers to a curve lying on a product of circles, not to a two-dimensional torus. It spans $\mathbb R^{2m}$, has a one-parameter simultaneous rotational symmetry, and is non-self-similar when $m\geq2$ because the ratios of its different coordinate radii change with time.

The rescaled curves $(-t)^{-1/2}\gamma_t$ tend to the radius-$\sqrt2$ circle with multiplicity $k_m$ as $t\to-\infty$, and with multiplicity one as $t\uparrow0$. Stryker--Sun prove
\[
 \sup_{t<0}\lambda(\gamma_t)=k_m\lambda(S^1)
 =k_m\sqrt{\frac{2\pi}{e}},
\]
which also bounds the entropy of every time slice. This is a supremum over time, not a formula for the entropy at each finite time. For the fixed center and Gaussian scale, direct integration gives
\[
 F_{0,1}(\gamma_t)=\sqrt\pi\left(\sum_j k_j^2r^{2k_j^2}\right)^{1/2}e^{t/2}.
\]
Indeed, $|\gamma_t|^2=-2t$ and $|\partial_\theta\gamma_t|^2=\sum_jk_j^2r^{2k_j^2}$. The centered Gaussian area is distinct from the supremum defining $\lambda$.

Source: \cite{CatSS19}, equation (2.1), Propositions 2.1--2.2, Corollary 2.3, and Theorem A.1, PDF pp.~3--5 and 8--9. The general source formula allows $k_1>1$; if the frequencies have common divisor $d>1$, that parametrization is a $d$-fold cover and its mass counts multiplicity. The embedded subfamily above avoids this ambiguity. The construction was already present in \cite{CatAAAW13}, Section 3, equations (11)--(14), PDF pp.~8--9, as acknowledged in \cite{CatSS19}, Remark 1.3. The interactive figure uses $(k_1,k_2)=(1,2)$ in $\mathbb R^4$. Its three-coordinate display is a projection, with time and fourth-coordinate rotation controls.
\end{example}

\chapter{Eternal shrinking helices}\label{sec:catalogue-eternal-helices}
\begin{example}[Helical curve-shortening flows]\label{ex:catalogue-eternal-helices}
Fix $m\geq1$ and integers $1\leq k_1<\cdots<k_m$. Let $r(t)>0$, $t\in\mathbb R$, solve the exact implicit equation
\[
 \log r(t)+\frac12\sum_{j=1}^m r(t)^{2k_j^2}=\frac m2-t.
\]
The choice of constant makes $r(0)=1$. Then
\[
 \Gamma_t(s)=
 \bigl(r^{k_1^2}\cos(k_1s),r^{k_1^2}\sin(k_1s),\ldots,
 r^{k_m^2}\cos(k_ms),r^{k_m^2}\sin(k_ms),s\bigr),
 \qquad s\in\mathbb R,
\]
is an eternal curve-shortening flow in $\mathbb R^{2m+1}$. In the simplest case $m=1$, $k_1=1$, it is the three-dimensional helix
\[
 \Gamma_t(s)=(r(t)\cos s,r(t)\sin s,s),\qquad
 \log r(t)+r(t)^2/2=1/2-t.
\]
These are exact parametrizations with an implicitly determined time scale, not numerical ODE profiles. The ordinary differential equation is
\[
 r'=-\frac{r}{1+\sum_j k_j^2r^{2k_j^2}}.
\]
The last coordinate makes every time slice a proper embedding of $\mathbb R$, and the induced metric is $(1+\sum_jk_j^2r^{2k_j^2})ds^2$, so it is complete. It is noncompact with two ends and codimension $2m$, and it spans the stated ambient space. Simultaneous rotations together with translation along the last coordinate give its screw symmetry. Its radius decreases to zero as $t\to+\infty$, giving local smooth convergence to the axis, while $r(t)\to\infty$ as $t\to-\infty$. At any fixed finite time the ends continue winding; they are not asymptotic to the axis as $|s|\to\infty$.

Proposition 2.7 of \cite{CatSS19} proves
\[
 \sup_{t\in\mathbb R}\lambda(\Gamma_t)=+\infty.
\]
This concerns the entire time family and does not say that every individual slice has infinite entropy. In fact, a direct calculation gives the exact fixed-time entropy below. Set $B=\sum_jr^{2k_j^2}$ and $D=\sum_jk_j^2r^{2k_j^2}$. Direct Gaussian integration, using $|\Gamma_t(s)|^2=B+s^2$, gives the exact centered value
\[
 F_{0,1}(\Gamma_t)=\sqrt{1+D}\,e^{-B/4};
 \qquad F_{0,1}(\Gamma_t)=\sqrt{1+r^2}\,e^{-r^2/4}
 \quad\hbox{when }m=k_1=1.
\]
For completeness, the same calculation at arbitrary Gaussian scale gives
\[
 \lambda(\Gamma_t)=\sqrt{1+D(t)}.
\]
Indeed, for a center $(a,z_0)\in\mathbb R^{2m}\times\mathbb R$, dropping the nonnegative squared distance in the first $2m$ coordinates bounds $F_{(a,z_0),\tau}$ by
$(4\pi\tau)^{-1/2}\sqrt{1+D}\int_{\mathbb R}e^{-(s-z_0)^2/(4\tau)}\,ds=\sqrt{1+D}$.
At the origin, $F_{0,\tau}=\sqrt{1+D}e^{-B/(4\tau)}$ tends to this upper bound as $\tau\to\infty$. Thus every finite-time helix has finite entropy, while the time supremum is infinite. This is a direct integration, additional to the cited time-supremum result.

Source: \cite{CatSS19}, equation (2.6), Propositions 2.6--2.7, and Theorem A.2, PDF pp.~6--7 and 9; see also \cite{CatAAAW13}, Section 3, PDF pp.~8--10. The interactive three-dimensional figure uses the simplest exact helix, with a time slider and the displayed range $-5\pi/2\leq s\leq5\pi/2$. The family is eternal even though a display uses only finitely many turns.
\end{example}

\chapter{Immortal annuli converging to a double plane}\label{sec:catalogue-double-plane-flow}
\begin{example}[Chen--Sun's noncompact multiplicity-two limit]\label{ex:catalogue-double-plane-flow}
Chen and Sun construct a nonstatic immortal mean curvature flow $M_t\subset\mathbb R^3$, $t\in[0,\infty)$, smooth for $t>0$, of connected embedded rotational surfaces. It is also symmetric across the limiting plane. Each positive-time slice consists of two increasing radial graphs joined smoothly at a neck of radius $a(t)>0$:
\[
 M_t=\{(\rho\cos\theta,\rho\sin\theta,\pm f(\rho,t)):
 \rho\geq a(t),\ \theta\in\mathbb R/2\pi\mathbb Z\}.
\]
This is a description of the constructed PDE solution, not an explicit formula for $f$. Away from the neck its upper graph satisfies
\[
 f_t=\frac{f_{\rho\rho}}{1+f_\rho^2}+\frac{f_\rho}{\rho}.
\]
Near the neck, the inverse graph $\rho=u(z,t)$ instead satisfies
\[
 u_t=\frac{u_{zz}}{1+u_z^2}-\frac1u,\qquad u_z(0,t)=0.
\]
The proof selects a critical initial surface by interpolation among rotational surfaces in finite cylinders and then takes a limit as the cylinder radius tends to infinity. Its initial trace is locally Lipschitz; choosing any positive time as the initial time gives an entirely smooth immortal flow. No closed formula, canonical parameter value, or uniqueness of the constructed critical flow is supplied.

The two graphical ends and the smooth connecting neck give topology $S^1\times\mathbb R$: genus zero and two ends. Properness follows from the graphs escaping to infinite radius; the smooth connected surface is complete. It remains in a fixed slab, and the two ends stay asymptotic to their original parallel planes. The neck radius tends to zero only as $t\to\infty$. On compact regions the flow converges to the middle plane with multiplicity two; the two sheets converge smoothly away from the origin. The convergence is not uniform over the entire noncompact surface. This is an infinite-time limit of the unrescaled flow, not a multiplicity-two finite-time tangent flow of a closed embedded surface.

Source: \cite{CatCS26R3}, Theorems 1.1 and 4.13, Section 1.1, equation (2.5), Proposition 4.6 and Remark 4.7, and Appendix A, arXiv v3 PDF pp.~1, 3, 6, 23--24, and 27--28. The topology and properness statements above are direct consequences of the source's two-graph parametrization. The paper gives no exact entropy or fixed-center Gaussian-area value for this constructed family. The interactive picture compares ordered finite-time stages and a separate infinite-time limit at a fixed viewing scale. A cutaway exposes the connecting neck; the compressed radial profile records the fixed far-end heights. The illustrated shapes and stage spacing are schematic, not computed critical PDE time slices or an asserted convergence rate.
\end{example}

\chapter{Immortal compact flows converging to a double sphere}\label{sec:catalogue-double-sphere-flow}
\begin{example}[Chen--Sun's product-manifold construction]\label{ex:catalogue-double-sphere-flow}
For each $n\geq2$, there is a smooth connected embedded immortal mean curvature flow of compact hypersurfaces in the product manifold
\[
 (S^n\times[-1,1],\ g_{S^n}+dy^2),
\]
where $S^n$ has unit round radius. The flow is defined for all $t\geq0$ and is nonstatic. Its hypersurface dimension is $n$ and its codimension in this ambient manifold is one. The construction stays in the smaller band $|y|\leq1/(2n)$, so the same flow is contained in $S^n\times\mathbb R$; it can also be placed in the closed product $S^n\times(\mathbb R/4\mathbb Z)$ without changing the metric near the flow. There is no free boundary on $y=\pm1$.

The solution is invariant under rotations of the sphere about its polar axis, reflection $y\mapsto-y$, and reflection across the sphere's equator. A closed profile $\sigma\mapsto(x(\sigma,t),y(\sigma,t))$ in $(0,\pi)\times\mathbb R$ generates it by
\[
 (\sigma,\omega)\longmapsto
 (\cos x(\sigma,t),\sin x(\sigma,t)\,\omega,y(\sigma,t)),
 \qquad \sigma\in S^1,\quad\omega\in S^{n-1}.
\]
The profile avoids both axes at every finite time, so this parametrization directly gives topology $S^1\times S^{n-1}$. Thus for $n=2$ the surfaces are genus-one tori with no ends. All slices are complete and proper by compactness. This is a topological deduction from the construction, not an assertion that the source states a separate topology theorem.

No closed time-dependent profile is supplied. Its upper half can be described using the solution of
\[
 v_t=\frac{v_{xx}}{1+v_x^2}+(n-1)\cot x\,v_x,
 \qquad v_x(\pi/2,t)=0,
\]
with an inverse graph near each neck, governed by
\[
 u_t=\frac{u_{yy}}{1+u_y^2}-(n-1)\cot u,
 \qquad u_y(0,t)=0.
\]
An interpolation argument selects an immortal solution. As $t\to\infty$, it converges to the stable minimal slice $S^n\times\{0\}$ with multiplicity two. Away from the north and south poles, the two sheets converge smoothly; the two necks collapse only at infinite time.

Source: \cite{CatCS25Man}, Theorems 1.1, 3.12, and 3.17, Section 2.1, Definition 3.10 and Lemma 3.11, and Section 3.3, arXiv v2 PDF pp.~1--4 and 11--15. Euclidean Gaussian area and Colding--Minicozzi entropy are not assigned to this intrinsic manifold flow; the cited paper gives no such invariant for these examples. In particular, its isometric embedding into $\mathbb R^{n+2}$ is not an unforced Euclidean mean curvature flow. The picture pairs three ordered finite-time stages with the infinite-time limit and their orbit-space profiles. In the $n=2$ drawing, the coordinate map $(p,y)\mapsto(1+2y)p$, for $p\in S^2$, displays the two sheets as inner and outer shells joined through both polar necks. This map is not an isometry: the radial separation illustrates the product coordinate, not a Euclidean MCF. A display cutaway exposes the inner sheet; its cut edges are not boundaries of the compact surface. The illustrated profiles are smooth schematic loops, not computed solutions or an asserted asymptotic formula. The stage slider assigns no physical times or decay rates. At every finite stage both neck circles remain positive, whereas the separate limit view shows one sphere with multiplicity two.
\end{example}

\chapter{Haslhofer--Hershkovits ancient ovals}\label{sec:catalogue-hh-ancient-ovals}
\begin{example}[Haslhofer--Hershkovits ancient ovals]\label{ex:catalogue-hh-ancient-ovals}
For $n\geq2$ and every $1\leq j\leq n-1$, Haslhofer--Hershkovits construct an ancient hypersurface flow $M_t\subset\mathbb R^{n+1}$, $-\infty<t<0$, with compact strictly convex embedded slices. Each slice is diffeomorphic to $S^n$, hence complete and proper with no ends. The construction has $O(j+1)\times O(n-j)$ symmetry and is uniformly Andrews noncollapsed. It is not self-similar and becomes round at extinction.

Under backward parabolic blow-down,
\[
 \lambda^{-1}M_{\lambda^2t}\longrightarrow S^j(\sqrt{2j|t|})\times\mathbb R^{n-j}\quad(\lambda\to\infty).
\]
The tip rescalings selected in Claim 3.5 of \cite{CatHH13} approach $\mathrm{Bowl}^{j+1}\times\mathbb R^{n-j-1}$, with the bowl normalized to unit tip curvature. Rescaled diameter is unbounded backward in time and curvature decay is type II. No closed formula or uniqueness follows from this construction. Gaussian area and entropy are finite for each compact slice; no fixed slice value is evaluated here.

Source: \cite{CatHH13}, Theorem 1.1 and Section 3. Huisken--Sinestrari's Theorem 3.2 in \cite{CatHS14} shows why the hypothesis $h_{ij}\geq\varepsilon H g_{ij}$ with one $\varepsilon>0$ for all ancient times cannot hold for these nonspherical ovals: a closed convex ancient flow with that uniform pinching is a shrinking sphere.

The illustration specializes to $n=2$, $j=1$. It numerically evolves the boundary of the Minkowski sum of a unit ball and an ellipsoid with semiaxes $(8,1/2,1/2)$, using rotational mean curvature flow. This is a finite smooth convex seed inspired by the long capped cylinders in Section 3, not a closed formula for the ancient oval. Extinction is shifted to $t=0$, and parabolic scaling makes the axial-to-equatorial radius ratio equal to $2$ at $t=-1$. The slider samples $-1000\leq t\leq-0.001$.

Original size shows contraction at a fixed spatial scale. The default view shows $(-t)^{-1/2}M_t$: the backward central limit is the cylinder of radius $\sqrt2$, while the extinction limit is the sphere of radius $2$. The right-tip view instead translates the tip to the origin and multiplies lengths by its mean curvature; the dashed comparison is a numerically integrated unit-speed bowl. These are different rescalings. Mesh refinement, time-tolerance refinement, and a shorter seed give consistent sampled profiles, but they do not certify an error to the ancient solution.
\end{example}

\chapter{Ancient flows from unstable minimal hypersurfaces}\label{sec:catalogue-minimal-ancient-flows}
\begin{example}[Ancient flows from unstable minimal hypersurfaces]\label{ex:catalogue-minimal-ancient-flows}
Let $M^n\subset\mathbb R^{n+1}$, $n\geq2$, be two-sided, properly embedded and minimal, unstable on a bounded domain, with bounded curvature and a uniform embedded tubular neighborhood. Assume additionally that it is asymptotically flat, periodic with compact fundamental domain, or periodic and asymptotically flat within its fundamental domain. Mramor--Payne produce two embedded ancient flows approaching $M$ smoothly and uniformly from opposite sides as $t\to-\infty$. Each has a sign of mean curvature, is nonconvex, and is not a soliton.

For a concrete three-dimensional example, use the classical Costa minimal surface. For sufficiently negative times the normal-graph construction retains its genus one and three ends. These slices are complete and proper; these conclusions follow from the bounded normal graph over the complete proper initial model. We record that early ancient interval $(-\infty,T)$, without asserting an eternal continuation for the Costa case. The theorem also applies to other Costa--Hoffman--Meeks surfaces and periodic minimal surfaces, but their differing topology is not collapsed into this case.

Source: \cite{CatMP19}, Theorem 1.1 and its following examples, pp.~3--4, and the normal-graph construction in Section 3. There is no closed time-dependent formula, uniqueness, or quantitative entropy value supplied. The distinct rotational eternal continuation from a catenoid is recorded separately in \ref{ex:catalogue-reapernoid}. The picture uses a numerical sampling of the classical Costa parametrization in \cite{CatMelko2010}, equation (17), with square lattice $\mathbb Z+i\mathbb Z$ and no additional spatial dilation. Its three ends are truncated for display.

To illustrate different early times, the viewer uses the linearization of mean curvature flow over the minimal surface. With $H=\operatorname{div}\nu$ and mean curvature vector $-H\nu$, a normal height satisfies
\[
 \partial_t u=(\Delta_M+|A|^2)u
\]
to first order. On a truncated Costa domain, let $\phi_D>0$ be the numerically computed first Dirichlet mode, normalized by $\max\phi_D=1$, with $(\Delta_M+|A|^2)\phi_D=\lambda_D\phi_D$ and $\lambda_D>0$. The displayed model is
\[
 X_{\mathrm{model}}^{\pm}(p,\tau)
 =X_M(p)\pm0.11e^\tau\phi_D(p)\nu_M(p),
 \qquad \tau=\lambda_D t\in[-5,0].
\]
The numerical value $\lambda_D\approx1.511$ depends on this scale, truncation and discretization; it is not a claimed exact eigenvalue of the complete Costa surface. The height decays backward and increases forward on either normal side. These are linearized model times, not computed nonlinear ancient slices or an asymptotic expansion proved for the flows in Theorem 1.1. The default display magnifies the normal displacement by three; the model-scale option removes that magnification. The endpoint $\tau=0$ is only a chosen small model amplitude, not a singular time. The two signed models illustrate the two sides; no reflection identity between the nonlinear flows is asserted.
\end{example}

\chapter{The eternal reapernoid}\label{sec:catalogue-reapernoid}
\begin{example}[The eternal reapernoid]\label{ex:catalogue-reapernoid}
For every $n\geq2$, Mramor--Payne construct an eternal embedded hypersurface flow $M_t\subset\mathbb R^{n+1}$, $t\in\mathbb R$, on the outside of a rotational minimal catenoid. It is $O(n)\times O(1)$ invariant, mean convex after a choice of orientation, nonconvex, and has uniformly bounded curvature. Backward in time it converges smoothly and uniformly to that catenoid; forward in time its distance from the rotation axis tends to infinity. It is not a translator or another soliton.

The rotational graphical construction gives complete proper noncompact slices with topology $S^{n-1}\times\mathbb R$ and two ends. For $n=2$ these are annuli in $\mathbb R^3$. In that dimension the pointed forward profile limit is a line and curvature tends to zero. For $n\geq3$, recentered meridians converge instead to a grim reaper with the same width as the higher-dimensional catenoid. These are pointed profile limits, not a claim that the unshifted surfaces converge globally to a fixed grim cylinder.

Source: \cite{CatMP19}, Theorem 1.3, pp.~4--5, and Section 4. The scale of the limiting catenoid is free; no explicit time-dependent profile, full uniqueness statement, or entropy value is recorded. The figure uses numerical solutions of the rotational evolution equation (2.7), rather than dilating a fixed catenoid.

Write the rotational hypersurface as
\[
 X(x,\omega,t)=(x,u(x,t)\omega),\qquad \omega\in S^{n-1},\qquad
 \partial_tu=\frac{u_{xx}}{1+u_x^2}-\frac{n-1}{u}.
\]
The reference catenoid has neck radius one. For $n=2$ its profile is exactly $w(x)=\cosh x$ on $\mathbb R$. Its outward parallel perturbation can be written parametrically as
\[
 x=s-\delta\tanh s,\qquad r=\cosh s+\delta\operatorname{sech}s,
 \qquad s\in\mathbb R,
\]
for sufficiently small $\delta>0$. These offsets are the initial surfaces for the approximating flows, not time slices of the limiting eternal flow. The numerical illustrations use $\delta=10^{-4}$, evolve forward, and shift time so that $u(0,0)=2$. Space is not rescaled from frame to frame. The surface views in $\mathbb R^3$ cover $t\in[-8,10]$ and the meridian views continue to $t=100$; the $n=3$ meridians in $\mathbb R^4$ cover $t\in[-2,40]$. These finite intervals and their time origin are plotting choices, not the maximal interval of the theorem.

The neck moves outward: $u_t>0$. In the convention $H=\operatorname{div}_M\nu$, the normal pointing away from the rotation axis has $H<0$ on moving slices, so the mean curvature vector $-H\nu$ points outward. The convexity of the meridian does not make the hypersurface convex: its meridional and rotational principal curvatures have opposite signs.

The late-time views subtract the neck radius, showing $u(x,t)-u(0,t)$ without dilation. For $n=2$ the limiting graph is zero. For the plotted $n=3$ case, set
\[
 a_3=\int_1^\infty\frac{dr}{\sqrt{r^4-1}},\qquad
 c=\frac{\pi}{2a_3},\qquad
 G(x)=-\frac{1}{c}\log\cos(cx),\quad |x|<a_3.
\]
Here $a_3$ is the coordinate half-width for the stationary equation with $w(0)=1$, and $G$ is the recentered grim-reaper limit. The dashed limit curves are distinct from the finite-time numerical profiles. Only $n=2$ and $n=3$ are computed; the existence and limiting-shape theorem applies to every $n\geq2$.

The computations use finite starting times, symmetry at the neck, and fixed values at artificial outer boundaries. Mesh, starting-offset, time-tolerance and domain comparisons are recorded with the figure data. They provide numerical evidence for these approximations, not a certified error estimate to an eternal solution.
\end{example}

\chapter{Ancient stacked pancakes}\label{sec:catalogue-stacked-pancakes}
\begin{example}[Ancient stacked pancakes]\label{ex:catalogue-stacked-pancakes}
For each $n\geq3$, Mramor--Yudowitz construct a smooth ancient flow of closed embedded hypersurfaces $M_t^n\subset\mathbb R^{n+1}$ for $t\in(-\infty,0]$. The slices are nonconvex, contained in a fixed slab, and $O(n)\times O(1)$ invariant. Their rotational radius profile is a graph over the axis with two local maxima separated by one local minimum. Topologically this is two lobes joined by one neck, giving $S^n$; compactness gives completeness, properness, and zero ends.

The endpoint $t=0$ is a regular chosen time, not an asserted extinction time. The theorem does not establish this construction for $n=2$, so the lowest ambient dimension is four. It also does not prove a precise backward decomposition into two separate ancient pancakes and a catenoid; the authors explicitly distinguish this motivation from their result. Existence uses a continuity and compactness argument. No closed formula, uniqueness, stability result, or uniform ancient entropy value is recorded. Each compact slice has finite entropy and Gaussian area.

Source: \cite{CatMY26}, Theorem 1.1, Figure 1, and Sections 2.2--3.1 of arXiv v3. To read the pictures, write a rotational slice as $(u(x,t)\omega,x)$, where $\omega\in S^{n-1}$ and $a(t)<x<b(t)$ is the axial coordinate. The positive radius $u$ has two maxima and a central neck minimum; it vanishes at the two poles. The first picture shows a schematic three-dimensional linear section, with a removable display cutaway exposing the neck. The second shows the meridian and a fixed containing slab. These original drawings illustrate the geometry and symmetry, not a numerical time slice or a mean curvature flow in $\mathbb R^3$; their proportions are not asserted by the theorem.
\end{example}

\chapter{Compact nonconvex ancient spiral curves}\label{sec:catalogue-ancient-spiral-curves}
\begin{example}[Compact nonconvex ancient spiral curves]\label{ex:catalogue-ancient-spiral-curves}
Zhang--Olson--Khan--Angenent construct compact embedded ancient curve-shortening flows $C_t\subset\mathbb R^2$, $t<0$, whose distant past follows a truncated rotating Yin--Yang spiral closed by a small grim-reaper-shaped cap. They are smooth simple closed curves, hence complete and proper, with no ends. They are nonconvex in the ancient regime and are neither shrinking circles nor ancient paperclips. No assertion of nonconvexity through a later maximal continuation is needed here.

The limiting rotating model has angular speed one. Its radial expansion $\mathcal R(\theta)=\sqrt{2\theta}+O(\theta^{-3/2})$ is asymptotic, not a formula for the constructed ancient curve. If $\Omega(t)$ is the paper's truncated spiral region and $\Lambda(t)$ is the enclosed ancient region, the Main Theorem proves
\[
 |\Lambda(t)\mathbin{\triangle}\Omega(t)|\leq C_\delta |t|^{-1+\delta}\qquad(t\to-\infty)
\]
for each $\delta>0$. This is an area estimate, not an exact pointwise parametrization. Gaussian area and entropy are finite at each compact slice; no uniform ancient bound is assigned.

Source: \cite{CatZOKA23}, Main Theorem and equations (1.1)--(1.3), pp.~1--3, and Appendix A. No uniqueness is claimed. The source also credits Charyyev's independent construction. The figure samples the asymptotic construction at $t=-24,-18,-12,-8$ on a common spatial scale. Its blue arms solve the rotating-soliton ODE numerically; the orange cap uses the corrected grim-reaper ansatz in Section 2. These are model snapshots, not a numerical solution of the ancient evolution equation or an illustration of extinction.

To read the motion, write $E(\theta)=(\cos\theta,\sin\theta)$ and $J(x,y)=(-y,x)$. The two arms are
\[
 Y_\pm(\theta,t)=\mathcal R(\theta-t\pm\pi/2)E(\theta),
 \qquad \theta-t\pm\pi/2>\pi/2,
\]
truncated near $\theta=-t$. Here $\mathcal R$ is the Appendix A solution of
\[
 \tfrac12\mathcal R^2-\theta+\arctan(\mathcal R'/\mathcal R)=0,
 \qquad \theta>\pi/2,
\]
with $\mathcal R\downarrow0$ as $\theta\downarrow\pi/2$. At a fixed point of the unrotated model the arms rotate counterclockwise with speed one. The cap center instead follows $R(t)E(-t)$, where $R(t)=\mathcal R(-2t)\sim2\sqrt{-t}$, so it moves clockwise and inward. Its leading width is $\pi/R(t)$. The inset uses $X=e^{-tJ}[(R(t),0)+R(t)^{-1}(u,v)]$; its dashed reference is the leading grim reaper
\[
 G(p)=(-\arcsin(\tanh p),-\log\cosh p),\qquad p\in\mathbb R.
\]
The orange curve includes the paper's normal correction $(-4t)^{-1}F(p)JG'(p)$, equations (2.19)--(2.21), and a smooth interpolation to the arms. These asymptotic formulas apply in the distant past; the figure's finite time interval carries no certified approximation-error bound.
\end{example}

\chapter{Eternal curves from chains of grim reapers}\label{sec:catalogue-grim-reaper-chains}
\begin{example}[Eternal curves from chains of grim reapers]\label{ex:catalogue-grim-reaper-chains}
Choose monotonically ordered heights $a_0<a_1<\cdots<a_m$ and horizontal shifts $C_1,\ldots,C_m$. Angenent--You construct embedded eternal planar curve-shortening flows which for sufficiently negative time are assembled asymptotically from alternating grim reapers on consecutive strips. To record a definite nontrivial subfamily, take $m=3$. With $v_k=\pi/(a_k-a_{k-1})$, the separate asymptotic models are
\[
 x=C_k+(-1)^k\left[-v_k^{-1}\log\sin(v_k(y-a_{k-1}))+v_kt\right],
 \qquad a_{k-1}<y<a_k.
\]
This formula describes the model translators, not the glued eternal solution.

The actual flow is a graph $x=u(y,t)$ on $a_0<y<a_3$, exists for every $t\in\mathbb R$, and is complete, properly embedded and noncompact, with topology $\mathbb R$ and two ends asymptotic to the outer horizontal lines. The two-end and properness statements follow from this graphical description and the outer asymptotes. Curvature changes sign in the nontrivial chain; total absolute curvature is uniformly bounded. There is no closed formula for $u$; the classification and flow entropy are qualified below.

Source: \cite{CatAY18}, equation (3), Section 3.1 and Lemmas 3.8--3.10, pp.~2 and 8--13. The paper also constructs immersed ancient chains; those are not being labeled embedded by this selected case. The illustration uses $a_k=k\pi$ and $C_k=0$, so that the three model tips have horizontal velocities $-1,+1,-1$. It shows a numerical finite-start approximation: smoothly joined reapers at $t=-10$ are evolved by curve shortening, and slices from $t=-6$ to $t=9$ are displayed at the same spatial scale. The outer ends are truncated far outside the viewing window. The joins follow the construction in Section 3.2 and Figure 10, using the orientation of equation (3) and Figure 9. This computation is not a closed formula or a certified approximation of the eternal limit; the later shapes are numerical observations, not an additional asymptotic theorem. In particular, $t=0$ is not an extinction time.

The classification in \cite{CatChoiSeoSuZhao2026}, Theorem~1.1, identifies graphical ancient trombones as the nonconvex embedded planar ancient flows with finite flow entropy, alongside the line, shrinking circle, paperclip and grim reaper. Here flow entropy means $\operatorname{Ent}[M]=\sup_t\lambda(M_t)<\infty$, not merely finite entropy of each slice. Corollary~1.5 records flow entropy $m+1$ for the $m$-reaper family; thus the selected $m=3$ case has flow entropy four. This does not identify the entropy or centered Gaussian area of an individual slice with four.
\end{example}

\chapter{Lawlor necks as stationary Lagrangian flows}\label{sec:catalogue-lawlor-necks}
Lawlor necks are the stationary endpoint of the Lagrangian neck construction.
\begin{example}[Stationary special Lagrangian necks]\label{ex:catalogue-lawlor-necks}
For $n\geq2$, $a_j>0$, and phases $\psi_j\in\mathbb R$, define
\[
 P_0(y)=\frac{\prod_{j=1}^n(1+a_jy^2)-1}{y^2},\quad
 P_0(0)=\sum_ja_j,\quad
 \phi_j(y)=\psi_j+\int_0^y\frac{dt}{(a_j^{-1}+t^2)\sqrt{P_0(t)}}.
\]
The map
\[
 X(x,y)=\left(x_j\sqrt{a_j^{-1}+y^2}e^{i\phi_j(y)}\right)_{j=1}^n,
 \qquad(x,y)\in\mathbb S^{n-1}\times\mathbb R,
\]
parametrizes a complete proper embedded special Lagrangian of dimension and codimension $n$. Its constant flow $L_t=X(\mathbb S^{n-1}\times\mathbb R)$ is eternal, $t\in\mathbb R$, since $\vec H=0$. This is the $\alpha=0$ case of \cite{CatJoyceLeeTsui2010}, Theorems~C--D, credited there to Lawlor.

It has two planar ends and topology $\mathbb S^{n-1}\times\mathbb R$; in $\mathbb C^2$ it is an annulus. The positive characteristic angles sum to $\pi$. They determine a neck up to scale, and equal $a_j$ give the rotational subfamily after a unitary change of coordinates. Calibration makes the neck area-minimizing for compactly supported competitors with the same boundary, hence stable for compactly supported variations.

For $n\geq3$, \cite{CatImagiJoyceOliveira2016}, Theorem~1.1, classifies closed-as-subsets, embedded, exact, asymptotically conical special Lagrangians with two transverse planar asymptotes and rate $\rho<0$, up to the specified special unitary rotation. The source's weaker rate $\rho<2$ additionally permits translation. These are conditional uniqueness statements. The drawing projects the exact $n=2$, $a_1=a_2=1$ map, whose common complex factor is $\sqrt{1+y^2/2}+iy/\sqrt2$.
\end{example}

\chapter{Immortal curve shortening with nonunique geodesic limits}\label{sec:catalogue-aryan-lee}
The ambient space here is a Riemannian two-sphere, not Euclidean space.
\begin{example}[A circle of limiting geodesics]\label{ex:catalogue-aryan-lee}
For every $C^\infty$ neighborhood of the round metric on $\mathbb S^2$, \cite{CatAryanLee2026}, Theorem~5.1, constructs a smooth antipodally invariant metric $g$ in that neighborhood and a smooth simple closed curve whose curve-shortening flow
\[
 \partial_t\gamma=\kappa_gN_g,\qquad t\in[0,\infty),
\]
stays embedded for all time. There is an embedded circle $Z\subset\mathbb {RP}^2$ and a smooth injective family $\Gamma:Z\to\mathcal E^\infty(\mathbb S^2)$ of simple closed $g$-geodesics, all of the same length, such that
\[
 \omega_{C^\infty}(\gamma)=\{\Gamma(z):z\in Z\}.
\]
Convergence is convergence of unparametrized embedded curves, or smooth convergence after reparametrization. In particular there is no unique limiting geodesic.

Each regular slice is a compact complete proper embedded circle, of dimension one and codimension one in $(\mathbb S^2,g)$, with no ends. The construction uses a prescribed length function and a spiraling slow mode, not a closed time-dependent formula. It does not assert rotational symmetry or stability of the evolving curve. The diagram is a schematic in parameter space; it is not a curve-shortening simulation on the round sphere. Euclidean Gaussian area and entropy are inapplicable to this intrinsic ambient flow.
\end{example}

\chapter{Lagrangian torus necks collapsing at infinite time}\label{sec:catalogue-su-tsai-wood}
These compact flows live in a flat complex torus, which is kept in the Riemannian-ambient group.
\begin{example}[Immortal Lagrangian connected sums]\label{ex:catalogue-su-tsai-wood}
Let $m\geq3$ and let $X_1,X_2$ be special Lagrangian $m$-tori of the same phase in the flat Calabi--Yau torus $\mathbb C^m/\Gamma$. Assume they intersect transversely at exactly one point, with tangent planes satisfying the Lawlor angle criterion (a type-one intersection). By \cite{CatSuTsaiWood2026}, the Main Theorem, Theorem~9.1, Proposition~9.5, and Corollary~9.6, a sufficiently small Lawlor-neck desingularization can be chosen to start a smooth embedded immortal Lagrangian flow $N_t$.

After translating the initial time, take $t\geq0$. Every finite-time slice is a compact complete proper embedded $m$-manifold of codimension $m$, with topology $\mathbb T^m\mathbin{\#}\mathbb T^m$ and no ends. As $t\to\infty$, the flow converges smoothly away from the intersection point to $X_1\cup X_2$. Its neck scale and curvature satisfy, for sufficiently large $t$,
\[
 c\,t^{-1/(m-2)}\leq\varepsilon(t)\leq C\,t^{-1/(m-2)},\qquad
 \max_{N_t}|A|=O(t^{1/(m-2)}).
\]
Rescaling about the intersection point by $\varepsilon(t)^{-1}$ gives the Lawlor neck locally smoothly, and the curvature becomes unbounded. These are asymptotic estimates from a parabolic gluing construction, not an exact formula for $N_t$. The singular limiting union does not make the regular slices nonembedded. No symmetry, uniqueness, or uniform mean-curvature decay is asserted. Euclidean Gaussian quantities are not assigned in this torus ambient space. The illustration is a neck-collapse schematic.
\end{example}

\chapter{Lagrangian sphere chains in Gibbons--Hawking spaces}\label{sec:catalogue-lee-tsai}
These are surface flows of codimension two in a hyperkahler four-manifold.
\begin{example}[Infinite-time curvature blow-up with vanishing mean curvature]\label{ex:catalogue-lee-tsai}
Use the Gibbons--Hawking setting of \cite{CatLeeTsai2026}, Section~2: the circle quotient contains the collinear centers $(a_j,0,0)$ with $-1=a_1<\cdots<a_N=1$, $N\geq3$, and
\[
 \varphi=\varphi_{\rm reg}+\frac12\sum_{j=1}^N
 ((x-a_j)^2+y^2+z^2)^{-1/2},\qquad
 g=\varphi^{-1}\eta^2+\varphi(dx^2+dy^2+dz^2),
 \quad d\eta=*d\varphi.
\]
The regular part is smooth, nonnegative, and harmonic on the source's quotient neighborhood. Let the initial Lagrangian sphere be the smooth circle lift of a positive concave graph $(x,u_0(x),0)$ joining the outer centers, with $u_0(\pm1)=0$. For fixed $\Lambda>0$, require the small-slope and curvature conditions of Theorem~6.2:
\[
 |u_0'|\leq c^{-1},\qquad -u_0''\leq\Lambda u_0\varphi(x,u_0,0),
\]
where $c$ is sufficiently large as supplied by that theorem. The reduced equation is
\[
 u_t=\frac{u_{xx}}{\varphi(x,u,0)(1+u_x^2)},\qquad u(\pm1,t)=0.
\]
The lifted flow exists smoothly for $t\geq0$, remains a compact embedded sphere, and converges to the $A_{N-1}$ chain of $N-1$ special Lagrangian spheres. Theorem~6.2 proves
\[
 \sup|\vec H|\longrightarrow0,\qquad
 0<\liminf_{t\to\infty}\frac{\log\max|A|}{\sqrt t}
 \leq\limsup_{t\to\infty}\frac{\log\max|A|}{\sqrt t}<\infty.
\]
Finite-time slices have genus zero, no ends, and are complete and proper. Junctions occur only in the infinite-time limit. The quotient drawing is schematic, not a numerical solution of this weighted equation. It is not ordinary planar curve shortening, and Euclidean entropy is inapplicable. The coauthor is Ping-Hung Lee, distinct from Tang-Kai Lee.
\end{example}

\chapter{Su--Wood stationary products of Lawlor profiles}\label{sec:catalogue-su-wood-stationary}
\begin{example}[Cohomogeneity-two special Lagrangians]\label{ex:catalogue-su-wood-stationary}
Let $m\geq3$, $d=m-1$, $I=(-\pi/(2d),\pi/(2d))$, $\mu_1,\mu_2>0$, and $\varphi\in\mathbb R$. Define
\[
 r(s)=(\cos(ds))^{-1/d},\quad \ell(s)=r(s)e^{is},\quad
 B(s)=\int_0^s r(v)^2\,dv.
\]
For $u\in S^{m-2}\subset\mathbb R^{m-1}$ and $(x,y)\in I^2$, put
\[
 X(u,x,y)=\left(e^{i\varphi}\mu_1\mu_2\ell(x)\ell(y)u,
 \frac{\mu_2^2r(y)^2-\mu_1^2r(x)^2}{2}
 +i\bigl(\mu_2^2B(y)-\mu_1^2B(x)\bigr)\right)\in\mathbb C^m.
\]
This is an exact integral parametrization of an embedded, exact special Lagrangian of dimension and codimension $m$. Thus $\vec H=0$ and $L_t=X(S^{m-2}\times I^2)$ for every $t\in\mathbb R$. It is the $G=SO(m-1)$ case of \cite{CatSuWood2026}, equation~(31), Theorem~5.6, Proposition~5.9 and Remark~5.10. It has cohomogeneity two under this group; for $m=3$ the topology is $S^1\times\mathbb R^2$, not a surface annulus.

The topology in general is $S^{m-2}\times\mathbb R^2$, with one end. The parametrization is proper: approaching either endpoint of $I$ makes $|\ell(x)\ell(y)|$ diverge, and the remaining factor $u$ is compact. The induced metric is therefore complete. The calibrated image is stable under compactly supported variations. No entropy value or global asymptotic classification is recorded here. The picture is the exact one-dimensional Lawlor profile used twice by the construction, not a projection of the full $m$-manifold.

The same paper constructs translators from a line and an Anciaux shrinker profile, and from an expander/shrinker pair. The line/expander case recovers the symmetric Joyce--Lee--Tsui translator of Example~\ref{ex:catalogue-jlt-translators}; it is not a second catalogue entry. For other symmetry groups some constructions have a one-dimensional singular set, so their smooth embeddedness must not be inferred from this $SO(m-1)$ case.
\end{example}

\chapter{Stolarski--Su finite-time Lagrangian neck pinches}\label{sec:catalogue-stolarski-su}
\begin{example}[Prescribed Type II blow-up rates]\label{ex:catalogue-stolarski-su}
For $n\geq4$ and every integer $K\geq2$, and for $n=3$ with $K\in\{2,3,4\}$, \cite{CatStolarskiSu2026}, Theorem~1.1 and Remark~1, construct smooth noncompact properly embedded exact almost-calibrated Lagrangian mean curvature flows $L(t)\subset\mathbb C^n$, $0\leq t<T<\infty$. Here we record the $SO(n)$-invariant subfamily, of dimension and codimension $n$. Its topology is $S^{n-1}\times\mathbb R$, with two ends, and every regular slice is complete. The topology follows from the global profile graph and its spherical group orbits.

The origin is a finite-time Type II singularity with
\[
 0<\liminf_{t\uparrow T}(T-t)^{K/2}\sup_{L(t)}|A|
 \leq\limsup_{t\uparrow T}(T-t)^{K/2}\sup_{L(t)}|A|<\infty.
\]
The unique tangent flow is a stationary multiplicity-one pair of transverse special Lagrangian planes. In contrast, the finer rescaling satisfies
\[
 (T-t)^{-K/2}L(t)\longrightarrow aL
 \quad\hbox{smoothly on compact sets},\qquad a>0,
\]
where $L$ is the rotational Lawlor neck. Up to a fixed unitary rotation, that limiting neck has the exact parametrization
\[
 X(u,\theta)=(\sin(n\theta))^{-1/n}e^{i\theta}u,
 \quad u\in S^{n-1},\quad0<\theta<\pi/n.
\]
This is a formula for the blow-up limit, not for $L(t)$. The actual flow is constructed by solving a nonlinear parabolic problem near a modulated neck of scale comparable to $(T-t)^{K/2}$, and has no closed formula. The initial Lagrangian angle oscillation can be arbitrarily small. These existence and convergence statements do not assert dynamical stability, uniqueness of the initial data, or a numerical entropy value. The illustration compares shrinking neck scales schematically; it is not a time integration of the PDE. The compact immersed variant mentioned in the source's Remark~3 is outside this entry.
\end{example}
